\chapter{Decision points DP1--DP7}
\label{ch:decisionpoints}

The architecture is organized around seven recurring decision points. Each has a natural Jev
primitive, some evidence it can be served, and a \term{risk class} that caps how much
authority the engine may have.

\begin{figure}[h]\centering
\small
\begin{tabular}{@{}llp{5.2cm}l@{}}
\toprule
ID & Decision & Primitive / evidence & Authority \\
\midrule
DP1 & intent / route & Choice; \IND{} within a few pts of a frontier LLM & may act \\
DP2 & \textbf{tool pre-selection} & Choice; \VENDOR{} skill-suggestion cookbook & narrows only \\
DP3 & action risk & Choice/Noul; \IND{} 91.7\% on 60 cases & \textbf{advisory only} \\
DP4 & \textbf{done / stuck} & Noul/Choice; evidence thin & control flow only \\
DP5 & observation triage & Noul per item; \VENDOR{} RAG cookbooks & filter only \\
DP6 & output verification & Choice/Noul; \IND{} agreement not accuracy & routes to review \\
DP7 & escalation & confidence from any DP & control flow \\
\bottomrule
\end{tabular}
\caption{The decision-point taxonomy. The MVP uses DP2 and DP4 (bold); DP3 is shadow-only,
the rest are deferred.\datalabel{Reported by docs/ARCHITECTURE.md \S3}}
\label{fig:dptaxonomy}
\end{figure}

\section{The design rule}

Jev is a candidate only for decisions whose answer space is \emph{closed}, whose relevant
state fits in about 32k tokens, and whose consequences can be bounded by code. Authority
scales inversely with risk.

\begin{figure}[h]\centering
\begin{tikzpicture}[node distance=5mm and 8mm, every node/.style={font=\scriptsize}]
  \node[frontend] (low) {\textbf{low risk}\\may act above\\threshold\\(DP1, DP7)};
  \node[infra, right=of low] (med) {\textbf{medium}\\narrows / filters /\\routes to review\\(DP2, DP5, DP6)};
  \node[security, right=of med] (high) {\textbf{high risk}\\advisory only;\\code + human decide\\(DP3, completion)};
  \draw[flow] (low)--(med); \draw[flow] (med)--(high);
  \node[below=8mm of med, font=\scriptsize\itshape, text width=8.5cm, align=center]
    {more authority to the engine \(\longleftarrow\)\hspace{1.5cm}\(\longrightarrow\) more authority to code + humans};
\end{tikzpicture}
\caption{Authority classes: the riskier the consequence, the less the engine may decide on
its own.\datalabel{Illustrative (docs/ARCHITECTURE.md \S3, \S10)}}
\label{fig:authority}
\end{figure}

\section{The two MVP decisions}

\begin{description}[leftmargin=1.4cm,style=nextline]
  \item[DP2 tool pre-selection] a Choice over the permission-filtered catalog, always with a
    ``\texttt{none\_of\_the\_above}'' escape hatch. High confidence exposes one tool; medium
    exposes the top-$k$; low exposes the full catalog. The LLM still writes the arguments. It
    is advisory: the LLM may pick outside the top-$k$, and the disagreement is logged.
  \item[DP4 done / stuck] a Noul/Choice over a projection of the task, acceptance criteria,
    and the last few actions. ``Done'' with high confidence only \emph{triggers} a
    deterministic completion verifier; it never ends the run by itself. v1 is scoped to
    ``done'' plus loop/repeat detection; ``drift from goal'' waits for an operational
    definition.
\end{description}

\begin{pitfall}
Options are permission-filtered per step, so the option \emph{set} changes every call.
``Fix the option order'' really means ``fix a stable sort key (by tool ID) that is invariant
under filtering.'' Otherwise calibration fitted on one option ordering will not transfer.
\end{pitfall}

\begin{checkbox}
Why does DP2 narrow options rather than pick the final tool, and why does DP4 never declare
success on its own? (Appendix~\ref{app:answers}.)
\end{checkbox}
